\section{封闭式评估（Close-ended Evaluation）}

\subsection{定义与典型任务}

封闭式任务是指\textbf{可能的答案数量有限}（通常少于 10 个），且通常只有一个或少数几个正确答案的任务。

\begin{figure}[H]
\centering
\includegraphics[width=0.95\textwidth]{slides-images/slide-004.jpg}
\caption{两种主要评估类型：封闭式评估（如情感分类，标签为 Negative/Positive）和开放式评估（如 GPT-2 生成的续写文本，无唯一正确答案）\protect\footnotemark}
\end{figure}
\footnotetext{Slides 第5页。}

\begin{importantbox}{封闭式任务 = 标准机器学习分类}
封闭式 NLP 任务在评估方法上与标准机器学习分类问题完全一致——可以使用准确率（Accuracy）、精确率（Precision）、召回率（Recall）、F1 分数、ROC 曲线、AUC 等。这并不意味着评估\textbf{简单}，只是说评估方法论没有 NLP 特有的特殊性。
\end{importantbox}

常见的封闭式 NLP 任务包括：

\begin{table}[H]
\centering
\begin{tabular}{lll}
\toprule
\textbf{任务} & \textbf{描述} & \textbf{典型基准} \\
\midrule
情感分析 & 判断文本的情感极性 & IMDb, SST \\
文本蕴含 & 假设是否被前提蕴含 & SNLI \\
词性标注 & 标注每个词的词性 & Penn Treebank \\
命名实体识别 & 识别文本中的实体 & CoNLL \\
共指消解 & 判断代词指代的对象 & WSC \\
问答 & 基于文本回答问题 & BoolQ, SQuAD \\
\bottomrule
\end{tabular}
\caption{常见封闭式 NLP 任务及其典型基准数据集}
\end{table}

\subsection{SuperGLUE：多任务综合基准}

SuperGLUE 是封闭式评估的``超级基准''，包含 8--9 个不同任务。评估方式是对各任务的性能取平均，得到一个综合排名。这曾是（约 2 年前之前）衡量模型通用语言能力的标准方式。

\begin{figure}[H]
\centering
\includegraphics[width=0.95\textwidth]{slides-images/slide-010.jpg}
\caption{SuperGLUE 基准中的任务示例：BoolQ（是/否问答）、RTE（文本蕴含）、WiC（词义消歧）、WSC（共指消解）等\protect\footnotemark}
\end{figure}
\footnotetext{Slides 第11页。}

\subsection{评估指标的选择至关重要}

\begin{warningbox}{不要盲目使用准确率}
以垃圾邮件分类为例：如果 90\% 的邮件是非垃圾邮件，那么一个始终预测``非垃圾''的分类器就能获得 90\% 的准确率——但它实际上什么也没分类。这就是为什么需要根据任务特点选择精确率、召回率或 F1 分数。
\end{warningbox}

多任务基准中的指标聚合同样需要谨慎：

\begin{itemize}
    \item 不同任务使用不同的指标（准确率、F1、相关系数等），直接取平均在数学上是不严谨的
    \item 讲者提到一个真实案例：某基准中有一列指标是``越低越好''，但人们在取平均时忘记加负号，直到很久之后才发现
    \item 不要假设已有的做法就是正确的——要独立思考评估方案的合理性
\end{itemize}

\subsection{虚假相关性（Spurious Correlations）}

\begin{figure}[H]
\centering
\includegraphics[width=0.95\textwidth]{slides-images/slide-012.jpg}
\caption{SNLI 中的虚假相关性：模型仅通过假设文本中是否包含否定词就能取得很好的性能，完全不需要看前提文本\protect\footnotemark}
\end{figure}
\footnotetext{Slides 第13页。}

2019 年的一篇论文发现，在 SNLI 文本蕴含任务中，\textbf{仅看假设文本}（忽略前提）就能取得很好的性能。原因在于：当标注者被要求写出``不被前提蕴含的假设''时，他们通常会加入否定词。因此模型学到了一个捷径——如果假设中有否定词，就预测``不蕴含''。

\begin{warningbox}{标签来源决定基准质量}
始终要追问：这些标签是怎么来的？标注过程中是否引入了系统性偏差？即使是被广泛使用的基准也可能存在严重的虚假相关性。不要因为``大家都在用''就认为一个基准是没有问题的。
\end{warningbox}

\subsection{本章小结}

封闭式评估是标准的机器学习分类问题，可以使用准确率、精确率、召回率等经典指标。然而，指标选择、指标聚合、标签质量和虚假相关性等问题仍然需要仔细考量。SuperGLUE 等多任务基准提供了综合评估的框架，但其简单平均的聚合方式并不完美。

%% ========== Part 3: 开放式评估 ==========

